\clearpage
\section{Lösungen}
\label{cuts:solutions}
Die Lösungen benutzen rationale Arithmetik, Ordnung und Mengen.

\subsection*{Zu Aufgabe 1: Schnittbedingungen prüfen}
\textbf{(a)} Nur $D$ erfüllt alle vier Bedingungen:
\begin{center}
\renewcommand{\arraystretch}{1.1}
\begin{tabular}{@{}lcccc@{}}
\toprule
Menge & Nicht leer & Echt & Abwärts & Ohne Maximum\\
\midrule
$A=\varnothing$ & nein & ja & ja & ja\\
$B=\mathbb Q$ & ja & nein & ja & ja\\
$C$ & ja & ja & ja & nein\\
$D$ & ja & ja & ja & ja\\
$E$ & ja & ja & nein & ja\\
\bottomrule
\end{tabular}
\end{center}
\textbf{(b)} $A$ ist leer, $B$ nicht echt und $1$ das Maximum von $C$.
Bei $E$ gilt $\frac12\in E$, aber $-1<\frac12$ und $-1\notin E$.
Zu $p\in D$ oder $p\in E$ liegt $(p+1)/2$ wieder in derselben
Menge und über $p$; in $\mathbb Q$ genügt $p+1$.

In $\varnothing$ kann kein Element die Abwärtsbedingung oder
die Bedingung zum fehlenden Maximum verletzen. Beide gelten
dort, doch die eigens verlangte Nichtleerheit fehlt.

\textbf{(c)} $C$ reicht nach links bis zum gefüllten Punkt $1$,
$D$ bis zum offenen Punkt $1$; $E$ liegt zwischen den offenen
Punkten $0$ und $1$. Ohne die vierte Bedingung wären $C$ und $D$
beide Beschreibungen desselben rationalen Randes.

\subsection*{Zu Aufgabe 2: Hauptschnitte}
\textbf{(a)} $q-1\in\widehat q$, aber $q\notin\widehat q$:
Der Hauptschnitt ist nicht leer und echt.
Aus $s<p<q$ folgt $s<q$, also Abwärtsabgeschlossenheit.
Für $p<q$ ist $r=(p+q)/2$ rational, und
\[
 r-p=q-r=\frac{q-p}{2}>0.
\]
Damit liegt $r$ in $\widehat q$ und über $p$; ein Maximum fehlt.

\textbf{(b)} Nur $1$ liegt in $\widehat{\frac32}$.
Ein größeres Element ist $\frac54$, denn $1<\frac54<\frac32$.

\textbf{(c)} Alle Elemente von $\widehat q$ sind kleiner als $q$,
aber $q$ gehört nicht dazu. Ein größtes Element muss zugleich
zur Menge gehören; eine obere Schranke muss das nicht.

\clearpage
\subsection*{Zu Aufgabe 3: Ordnung und Einbettung}
\textbf{(a)} Es gilt
$\widehat{-\frac12}\subsetneq\widehat{\frac23}$.
Zum Beispiel gehört $0$ nur zum rechten Schnitt.

\textbf{(b)} Ist $p\le q$ und $x\in\widehat p$, dann
$x<p\le q$. Also liegt $x$ in $\widehat q$.
Umgekehrt sei $\widehat p\subseteq\widehat q$.
Bei $p>q$ läge $q$ in $\widehat p$, aber nicht in $\widehat q$.
Das widerspricht der Inklusion. Somit gilt $p\le q$.

\textbf{(c)} Für $p<q$ liefert Teil~(b) die Inklusion;
$p\in\widehat q\setminus\widehat p$ macht sie echt.
Aus echter Inklusion folgt nach Teil~(b) $p\le q$;
$p=q$ würde gleiche Hauptschnitte liefern und ist ausgeschlossen.

Gleiche Hauptschnitte liefern beide Inklusionen, also
$p\le q$ und $q\le p$, folglich $p=q$.
Die Umkehrung folgt direkt aus der Definition.
Damit ist die Abbildung $q\mapsto\widehat q$ injektiv
und erhält und erkennt die Ordnung. Genau dies leistet
hier die Einbettung.

\subsection*{Zu Aufgabe 4: Ein Zeuge genügt}
\textbf{(a)} Sei $a\in A\setminus B$ und $b\in B$.
Bei $b=a$ läge $a$ in $B$.
Bei $a<b$ läge $a$ wegen der Abwärtsabgeschlossenheit von $B$
ebenfalls in $B$. Beide Fälle sind unmöglich.
Die totale rationale Ordnung lässt nur $b<a$ übrig.

\textbf{(b)} Aus $b<a$ und $a\in A$ folgt $b\in A$.
Da $b\in B$ beliebig war, ist $B\subseteq A$.
Der Zeuge $a\in A\setminus B$ zeigt sogar $B\subsetneq A$.

\textbf{(c)} Entweder gilt bereits $A\subseteq B$,
oder es gibt ein $a\in A\setminus B$.
Im zweiten Fall liefert Teil~(b) $B\subseteq A$.
Damit sind je zwei Schnitte vergleichbar.

Für beliebige Teilmengen gilt das nicht:
Weder $\{0\}\subseteq\{1\}$ noch $\{1\}\subseteq\{0\}$.
Diese Mengen sind nicht abwärts abgeschlossen;
der entscheidende Schluss aus Teil~(a) steht dort nicht zur Verfügung.

\clearpage
\subsection*{Zu Aufgabe 5: Zwei Schnitte vereinigen}
\textbf{(a)} Nach Aufgabe~3 ist $A\subsetneq B$.
Daher ist $A\cup B=B=\widehat{\frac23}$.
Die Vereinigung hat dieselbe Skizze wie $B$ und ist
die kleinste obere Schranke von $\{A,B\}$.

\textbf{(b)} Nach Aufgabe~4 gilt $A\subseteq B$ oder
$B\subseteq A$. Im ersten Fall ist $A\cup B=B$,
im zweiten $A\cup B=A$. Die Vereinigung ist also ein Schnitt.

\textbf{(c)} Beide Mengen liegen in ihrer Vereinigung.
Ist $V$ eine obere Schranke, so enthält $V$ sowohl $A$ als auch $B$.
Jedes Element von $A\cup B$ liegt in einer dieser Mengen und damit
in $V$. Also ist $A\cup B\subseteq V$.
Dies beweist die Supremumseigenschaft.
Da die Vereinigung hier selbst eines der beiden Familienglieder
ist, ist sie zugleich das Maximum der Familie.

\subsection*{Zu Aufgabe 6: Supremum ohne Maximum}
\textbf{(a)} $\widehat0$ gehört zu $\mathcal A$.
Für jedes $r<1$ gilt nach Aufgabe~3
$\widehat r\subseteq\widehat1$.
Somit ist die Familie nicht leer und nach oben beschränkt.

\textbf{(b)} Liegt $p$ in der Vereinigung, so gibt es ein
$r<1$ mit $p\in\widehat r$. Also ist $p<r<1$ und
$p\in\widehat1$.
Umgekehrt sei $p\in\widehat1$, also $p<1$.
Für $r=(p+1)/2$ gilt $p<r<1$.
Damit gehört $\widehat r$ zur Familie und $p$ zu diesem Schnitt.
Also liegt $p$ in der Vereinigung. Insgesamt:
\[
 \bigcup\mathcal A=\widehat1.
\]

\textbf{(c)} $\widehat1$ ist bereits eine obere Schranke.
Jede weitere obere Schranke $V$ enthält alle Familienglieder,
also auch deren Vereinigung $\widehat1$.
Damit ist $\widehat1$ die kleinste obere Schranke.

\textbf{(d)} Aus $\widehat1=\widehat r$ würde nach Aufgabe~3
$r=1$ folgen; in der Familie ist jedoch $r<1$ vorgeschrieben.
Zu jedem $r<1$ liefert $t=(r+1)/2$ ein $r<t<1$ und damit
\[
 \widehat r\subsetneq\widehat t\in\mathcal A.
\]
Kein Familienglied ist das größte. Das Supremum existiert
und liegt außerhalb der Familie.

\clearpage
\subsection*{Zu Aufgabe 7: Leere und unbeschränkte Familien}
\textbf{(a)} Jedes Familienglied ist eine Teilmenge von $\mathbb Q$;
also ist auch die Vereinigung darin enthalten.
Für jedes $p\in\mathbb Q$ gilt aber
\[
 p\in\widehat{p+1}\in\mathcal H.
\]
Somit enthält die Vereinigung jede rationale Zahl und ist ganz
$\mathbb Q$. Sie ist kein Schnitt, weil sie nicht echt ist.

\textbf{(b)} Eine obere Schranke $U$ müsste jeden Hauptschnitt
und damit jedes $p\in\widehat{p+1}$ enthalten.
Daraus folgte $U=\mathbb Q$, im Widerspruch zur Echtheit von $U$.
Die Familie hat also keine reelle obere Schranke.

\textbf{(c)} Die Vereinigung der leeren Familie ist $\varnothing$:
Kein Element kann aus einem Familienglied stammen.
Sie ist wegen fehlender Nichtleerheit kein Schnitt.
Jeder Schnitt $U$ ist trotzdem eine obere Schranke der leeren
Familie. Die Bedingung, dass jedes Familienglied in $U$ liegt,
hat hier keinen einzigen zu prüfenden Fall.

\textbf{(d)} Sei $U$ ein beliebiger Schnitt.
Wegen seiner Nichtleerheit gibt es ein $u\in U$.
Die Abwärtsabgeschlossenheit liefert
\[
 \widehat u\subseteq U.
\]
Wegen $u\in U$ und $u\notin\widehat u$ ist diese Inklusion echt.
Zu jedem Schnitt gibt es also einen kleineren Schnitt.

Da alle Schnitte obere Schranken der leeren Familie sind,
kann keiner die kleinste obere Schranke sein:
$\widehat u$ ist stets eine noch kleinere obere Schranke als $U$.
Die leere Familie besitzt kein reelles Supremum.

\textbf{Merksatz.} Jede \emph{nichtleere, nach oben beschränkte}
Familie reeller Schnitte besitzt ein reelles Supremum.
Ohne Beschränktheit kann die Vereinigung ganz $\mathbb Q$ werden;
ohne Nichtleerheit wird sie leer.
Beide Ausschlüsse sind Bedingungen an die Familie.

\clearpage
\subsection*{Zu Aufgabe 8: Vollständigkeit}
\textbf{(a)} Wir prüfen die Schnittbedingungen für
$S=\bigcup\mathcal A$.
\begin{enumerate}[label=\arabic*.]
 \item Alle Familienglieder bestehen aus rationalen Zahlen.
 Daher gilt $S\subseteq\mathbb Q$.
 \item Weil die Familie nicht leer ist, gibt es ein
 $A_0\in\mathcal A$. Weil $A_0$ ein Schnitt ist,
 gibt es ein $a\in A_0$. Dann liegt $a$ auch in $S$.
 \item Alle Familienglieder liegen in $U$; daher $S\subseteq U$.
 Weil $U\ne\mathbb Q$ ist, kann auch $S$ nicht ganz $\mathbb Q$ sein.
 \item Sei $p\in S$ und $q\in\mathbb Q$ mit $q<p$.
 Dann liegt $p$ in einem $A\in\mathcal A$.
 Dessen Abwärtsabgeschlossenheit liefert $q\in A\subseteq S$.
 \item Sei $p\in S$. Wähle wieder ein Familienglied $A$ mit
 $p\in A$. Da $A$ kein größtes Element hat, gibt es
 ein $r\in A$ mit $p<r$. Dieses $r$ liegt auch in $S$.
\end{enumerate}
Damit ist $S$ tatsächlich ein reeller Schnitt.

\textbf{(b)} Jedes Familienglied ist in seiner Vereinigung
enthalten. Also ist $S$ eine obere Schranke.
Sei nun $V$ eine beliebige obere Schranke und $p\in S$.
Dann gehört $p$ zu einem $A\in\mathcal A$.
Aus $A\subseteq V$ folgt $p\in V$.
Somit ist $S\subseteq V$, also
\[
 \boxed{\sup\mathcal A=\bigcup\mathcal A.}
\]
Der Beweis zeigt sowohl, dass die rechte Seite ein zulässiges
reelles Objekt ist, als auch, dass sie die gewünschte Ordnungsrolle hat.

\textbf{(c)} Die Nichtleerheit der Familie wird in Schritt~2
benutzt; die obere Schranke $U$ in Schritt~3.
Die übrigen Eigenschaften gehen jeweils von einem Familienglied
auf seine Vereinigung über.

Es wird keine Folge erzeugt und kein Grenzwert vorausgesetzt.
Die Vereinigung ist eine direkt definierte Menge.
Inklusion erklärt zugleich ihre Lage über allen Familiengliedern
und unter jeder anderen oberen Schranke.
